%==============================================================================
% jarvis-boxes.tex
% Admonition boxes for the Jarvis-HEP manual, built on tcolorbox.
%
% Four flavours, used sparingly (see the style guide in the writing plan):
%   \begin{jnote} ... \end{jnote}        general clarification
%   \begin{jtip} ... \end{jtip}          best-practice advice
%   \begin{jwarn} ... \end{jwarn}        foot-guns / things that bite
%   \begin{jschema}{File_schema.json} ... \end{jschema}
%                                        "verify these keys against the schema"
%
% Boxes are sized to the Tufte text column.  For a full-width box, wrap the
% environment in {fullwidth}.
%==============================================================================

\usepackage[most]{tcolorbox}
\usepackage{soul}   % \hl{}: line-breakable background highlight, for inline \cli{}

% Reuse the palette defined in jarvis-listings.tex.
\providecolor{jNote}{HTML}{1F4E79}
\providecolor{jTip}{HTML}{2E7D32}
\providecolor{jWarn}{HTML}{B3261E}
\providecolor{jSchema}{HTML}{6A1B9A}
\providecolor{jOpt}{HTML}{B8860B}   % \opt badge -- goldenrod, legible as text/border/fill

\tcbset{
  jbox/.style={
    enhanced, breakable,
    boxrule=0pt, leftrule=2.5pt,
    arc=1pt, outer arc=1pt,
    left=8pt, right=6pt, top=5pt, bottom=5pt,
    fonttitle=\sffamily\bfseries\footnotesize,
    fontupper=\small,
    before skip=0.8\baselineskip, after skip=0.8\baselineskip,
  }
}

% Title bar: a LIGHT tint of the accent colour, with the accent itself as the
% (dark) title text -- so the heading reads but is no longer a heavy dark band.
\newtcolorbox{jnote}{%
  jbox, colback=jNote!4, colframe=jNote,
  colbacktitle=jNote!15, coltitle=jNote!90!black,
  title={Note}}

\newtcolorbox{jtip}{%
  jbox, colback=jTip!5, colframe=jTip,
  colbacktitle=jTip!15, coltitle=jTip!90!black,
  title={Tip}}

\newtcolorbox{jwarn}{%
  jbox, colback=jWarn!5, colframe=jWarn,
  colbacktitle=jWarn!15, coltitle=jWarn!90!black,
  title={Warning}}

% Schema-check box takes the schema file name as a mandatory argument.
\newtcolorbox{jschema}[1]{%
  jbox, colback=jSchema!4, colframe=jSchema,
  colbacktitle=jSchema!15, coltitle=jSchema!90!black,
  title={Verify against \texttt{#1}}}

%==============================================================================
% Shell-command box: a terminal-styled wrapper for CLI snippets, matching the
% call-out boxes above (light title bar + accent left rule + light body).
% Uses tcblisting (listings engine) with the `shellinner' style from
% jarvis-listings.tex; \faTerminal is provided by fontawesome5 (loaded later but
% expanded only at use time).  Use \begin{shellbox} ... \end{shellbox}.
%==============================================================================
\definecolor{jShell}{HTML}{334155}   % terminal slate

% \cli{}'s real definition (jarvis-preamble.tex has a placeholder, overridden
% here once jShell/soul exist): a shell fragment mentioned mid-sentence gets the
% same slate identity as a shellbox, as a highlighted "chip" -- \hl{} (not
% \colorbox/\tcbox) specifically because it stays breakable at spaces, and real
% commands run up to ~45 characters (e.g. \cli{pip install -e
% '.[distributed,plot,dev]'}), too long to risk as one unbreakable box in the
% narrow tufte column. \color{jShell} is a plain declaration, not \textcolor{}{} --
% soul's \hl reconstructs its argument token-by-token and chokes on a nested
% braced-argument macro like \textcolor{X}{...} inside it (confirmed by trying).
\sethlcolor{jShell!10}
\renewcommand{\cli}[1]{{\color{jShell}\ttfamily\hl{#1}}}

%==============================================================================
% YAML "window" box: a GUI-editor look for YAML listings -- a title bar with
% traffic-light dots and a filename, faint indentation guides, and light line
% numbers.  Use \begin{yamlwin}[filename] ... \end{yamlwin} (the optional title
% defaults to "YAML").  The yamlwindow style (jarvis-listings.tex) colours
% expression values pink and &J paths dark green.
%==============================================================================
\definecolor{jWinFrame}{HTML}{CDD2D8}  % window border
\definecolor{jWinBar}{HTML}{ECEEF1}    % title-bar background
\definecolor{jWinTitle}{HTML}{5A6470}  % title-bar text
\definecolor{tlRed}{HTML}{FF5F56}
\definecolor{tlYel}{HTML}{FFBD2E}
\definecolor{tlGrn}{HTML}{27C93F}
\newcommand{\trafficlights}{\tikz[baseline=-0.6ex]{%
  \fill[tlRed](0,0)circle(2.1pt);\fill[tlYel](7pt,0)circle(2.1pt);\fill[tlGrn](14pt,0)circle(2.1pt);}}

\newtcblisting{yamlwin}[1][YAML]{%
  enhanced, breakable,
  colback=jBg, colframe=jWinFrame,
  % `colupper' reasserts the key colour as the ambient text colour on every
  % page a listing is split across -- without it, `breakable' loses the
  % \color set inside listings' `basicstyle' the moment the box crosses a
  % page boundary, and the continuation renders in plain black. Text-level
  % overrides (string/keyword/expression/path colours) are unaffected either
  % way, since those are applied token-by-token as each line is typeset.
  colupper=jYKey,
  boxrule=0.7pt, arc=4pt, outer arc=4pt,
  boxsep=0pt, left=4pt, right=5pt, top=5pt, bottom=5pt,
  colbacktitle=jWinBar, coltitle=jWinTitle,
  fonttitle=\sffamily\footnotesize,
  toptitle=4pt, bottomtitle=4pt,
  title={\trafficlights\hspace{0.8em}{\ttfamily\footnotesize #1}},
  before skip=0.9\baselineskip, after skip=0.9\baselineskip,
  % faint vertical indentation guides (one per 2-space level), like an editor
  underlay={\begin{tcbclipinterior}\foreach \k in {1,...,7}{%
    \draw[jGuide,line width=0.4pt]
      ([xshift=2.4em+\k*1em]interior.north west)--([xshift=2.4em+\k*1em]interior.south west);}%
    \end{tcbclipinterior}},
  listing only,
  listing options={style=yamlwindow},
}

% RETIRED 2026-07-25 in favour of `terminal' + `jcmd' below, which unifies the three
% shapes a transcript takes (command only / output only / command then output) under
% one chrome. Kept defined so an old draft still compiles; do not introduce new uses.
\newtcblisting{shellbox}{%
  enhanced, breakable,
  boxrule=0pt, leftrule=2.5pt, arc=1pt, outer arc=1pt,
  boxsep=0pt, left=8pt, right=6pt, top=4pt, bottom=4pt,
  colback=jShell!4, colframe=jShell,
  colbacktitle=jShell!15, coltitle=jShell!90!black,
  fonttitle=\sffamily\bfseries\footnotesize,
  toptitle=3.5pt, bottomtitle=3.5pt,
  title={\raisebox{-0.2pt}{\scalebox{0.8}{\faTerminal}}\ \ Shell},
  before skip=0.8\baselineskip, after skip=0.8\baselineskip,
  listing only,
  listing options={style=shellinner},
}

%%
% --- The unified terminal box -------------------------------------------------
% ONE chrome for all three shapes a terminal transcript takes in this book. The
% box supplies the frame, the title bar and (when both halves are present) the
% divider; the two inner listing environments supply the only thing that differs,
% namely whether a line is something you TYPE or something a program PRINTED.
%
%   commands only          \begin{terminal} \begin{jcmd} ... \end{jcmd} \end{terminal}
%   output only            \begin{terminal} \begin{jout} ... \end{jout} \end{terminal}
%   commands then output   \begin{terminal} \begin{jcmd} ... \end{jcmd}
%                            \termsep      \begin{jout} ... \end{jout} \end{terminal}
%
% The optional argument renames the title bar (default `Terminal'); use it when a
% block is a named artifact rather than a session, e.g. \begin{terminal}[Jarvis2 -h].
%
% Why a plain tcolorbox hosting two lstlistings, and not a tcblisting: a
% tcblisting writes its whole body out verbatim, so a \tcblower inside one is
% printed as the literal text "\tcblower" instead of splitting the box (verified).
% A plain tcolorbox scans its body normally, which lets each half run through its
% own listings style AND lets tcolorbox draw its native segmentation line between
% them -- that faint dashed rule is the divider, and it costs nothing.
\lstnewenvironment{jcmd}{\lstset{style=shellinner}}{}
\lstnewenvironment{jout}{\lstset{style=termout}}{}
\newcommand{\termsep}{\tcblower}

\newtcolorbox{terminal}[1][Terminal]{%
  enhanced, breakable,
  boxrule=0pt, leftrule=2.5pt, arc=1pt, outer arc=1pt,
  boxsep=0pt, left=8pt, right=6pt, top=4pt, bottom=4pt,
  colback=jShell!4, colframe=jShell,
  colbacktitle=jShell!15, coltitle=jShell!90!black,
  fonttitle=\sffamily\bfseries\footnotesize,
  toptitle=3.5pt, bottomtitle=3.5pt,
  title={\raisebox{-0.2pt}{\scalebox{0.8}{\faTerminal}}\ \ #1},
  before skip=0.8\baselineskip, after skip=0.8\baselineskip,
  % --- the command/output divide ---
  % The two halves are told apart by ink, not by background: commands in jShell,
  % output in the lighter jOut (see `termout' in jarvis-listings.tex), with this
  % faint dashed rule between them. A background tint for the lower half was tried
  % and rejected -- colbacklower needs skin=bicolor, and bicolor strokes the whole
  % lower fill path with `segmentation style', which drags the dashes around the
  % box's bottom edge as well.
  % `middle' is the total vertical room the divider occupies (the rule sits in the
  % centre of it), so it is the only knob for the divider's height. Compared 6/3/1.5/0pt
  % side by side: 0pt makes the rule read as an underline of the last command rather
  % than a separator, 6pt was airier than the transcript needs. 2pt keeps a hairline of
  % air on both sides while costing almost nothing vertically.
  middle=2pt,
  lower separated=true,
  segmentation style={draw=jShell!45, line width=0.4pt,
                      dash pattern=on 1.4pt off 1.6pt},
}
